\scanpage{78}

% L01 -- top edge
\[
\widetilde{s}_0=s_0-\unsure{o|\sigma|c}
\]

% L02 -- upper left, crossed diagonally
\deleted{\[
\begin{aligned}
s'_0&=\rho\widetilde{s}_0,\\
f_V(s'_0)&
\end{aligned}
\]}

% L03 -- top, evaluation
\[
f_V(\widetilde{s}_0)=\deleted{\unsure{f_{X}}}f(s_0-o)=
 \underbrace{f(s_0)}_{\substack{\Vert\\0}}+\underbrace{f(o)}_{-(\sigma'\rho')/\rho^2}-
 \underbrace{\varphi(s_0,o)}_{-\sigma/\rho}
=-\frac{\sigma'\rho'}{\rho^2}+\frac{\sigma}{\rho}
\]

% L04 -- upper third, two parity rows
\[
\left\{\begin{aligned}
n\text{ pair}\quad&\rho=2\rho',\ \sigma^{\deleted{\unsure{2}}}=\sigma'\\
n\text{ impair}\quad&\rho=\rho',\ \sigma=2\sigma'
\end{aligned}\right.\quad
\left\{\begin{aligned}
f_V(\widetilde{s}_0)&=-\tfrac14\tfrac{\sigma'}{\rho'}+\tfrac12\tfrac{\sigma'}{\rho'}=\tfrac14\tfrac{\sigma'}{\rho'}\\
f_V(\widetilde{s}_0)&=-\tfrac{\sigma'}{\rho'}+2\tfrac{\sigma'}{\rho'}=\tfrac{\sigma'}{\rho'}
\end{aligned}\right.
\]

% L05 -- upper middle
\[
\begin{aligned}
\text{Soit}\quad s'_0&=-\rho\widetilde{s}_0=\deleted{\unsure{-}}
(\underbrace{\zeta_0}_{\substack{\Vert\\\sigma}}a_0+\zeta_1a_1+\cdots+\underbrace{\zeta_n}_{\substack{\Vert\\1}}a_n)\\
f_V(s'_0)&=\rho\sigma-\rho'\sigma'=\rho'\sigma'
\end{aligned}
\]

% L06 -- upper middle, box at left
\[
\boxed{f_V(s'_0)=\rho'\sigma'}
\]

% L07 -- middle, right of box
\[
\varphi_V(\widetilde{s}_0,a_i)=\varphi_{\deleted{\unsure{V}}}(s_0-o,a_i)=
\underbrace{\varphi(s_0,a_i)}_{\substack{0\text{ si }i\ne0\\-1\text{ si }i=0}}
-\underbrace{\varphi(o,a_i)}_{0\ \forall i}
=\begin{cases}-1&i=0\\0&1\le\unsure{j|i}\le n\end{cases}
\]

% L08 -- middle left, two pairs
\[
\left\{\begin{aligned}\widetilde{s}_i&=s_i-o\\s'_i&=-\rho\widetilde{s}_i\end{aligned}\right.
\qquad\left\{\begin{aligned}
\widetilde{s}_i&=\tau_{i-1}^V\widetilde{s}_{i-1}\\
s'_i&=\tau_{i-1}^V s'_{i-1}
\end{aligned}\right.\qquad(1\le i\le n)
\]

% L09 -- middle right, bracket
\[
\begin{aligned}
\varphi_V(s'_0,a_i)&=\begin{cases}\rho&\unsure{j|i}=0\\0&1\le i\le n\end{cases}\\
\langle\widetilde{s}_0,a_i^\vee\rangle&=\langle\widetilde{s}_0,a'_i\rangle=\\
\langle s_0-o,a'_i\rangle&=\langle s_0,a'_i\rangle\\
&=\begin{cases}1&i=0\\0&1\le i\le n\end{cases}
\end{aligned}
\]

% L10 -- middle left, sums
\[
\left\{\begin{aligned}
\widetilde{s}_i&=(s_i-s_0)+(s_0-o)=\widetilde{s}_0+(a_0+\cdots+a_{i-1})\\
s'_i&=s'_0\unsure{-}\rho(a_0+\cdots+a_{i-1})
\end{aligned}\right.
\]

% L11 -- lower middle, first edge
\[
\begin{aligned}
s'_1&=s'_0-\rho a_0\\
\deleted{\unsure{f}}s'_0+s'_1&=2s'_0-\rho a_0\\
f_V(s'_0+s'_1)&=4\underbrace{f_V(s'_0)}_{\rho'\sigma'}+\rho^2
\underbrace{f_V(a_0)}_1\unsure{-}2\rho\underbrace{\varphi_V(s'_0,a_0)}_\rho
=4\rho'\sigma'-\rho^2=
\end{aligned}
\]

% L12 -- below first edge, parity split
\[
f_V(s'_0+s'_1)=\begin{cases}
n\text{ pair}:&4\rho'(\sigma'-\rho')\\
n\text{ impair}:&\rho'(4\sigma'-\rho')
\end{cases}
\]

% L13 -- lower third
Donc Pour que l'homothétie de rapport $\lambda$ transforme la sphère $\Sigma$ circonscrite aux $s'_i$ en celle circonscrite aux \unsure{transformés} de $s'_0+s'_1$, il f. et s. que

% L14 -- lower third, final ratios
\[
\lambda^2\rho'\sigma'=\deleted{\rho'(4\sigma'-\rho')}
\left\{\begin{aligned}
&4\rho'(\sigma'-\rho')\quad\text{i.e.}\quad
\lambda=2\sqrt{1-\rho'/\sigma'}\quad n\text{ pair}\\
&\deleted{\text{i.e. }\lambda=}\\
&\rho'(4\sigma'-\rho')\qquad\lambda=2\sqrt{1-\rho'/4\sigma'}\quad n\text{ imp}
\end{aligned}\right.
\]
% archive sheet number: 77
